\documentclass[12pt]{article}
\usepackage{amsmath, amssymb}

\begin{document}

\title{Lecture Scribe\\
CSE400 – Fundamentals of Probability in Computing\\
Lecture 18: Marginal PMF, Correlation, Covariance, Joint Gaussian Random Variables}
\author{Instructor: Dhaval Patel, PhD}
\date{March 12, 2026}
\maketitle

\section{Marginal PMF}

\subsection{Definition}

For discrete random variables $X$ and $Y$, the joint PMF is:
\[
p_{XY}(x,y) = P(X = x, Y = y)
\]

The marginal PMFs are obtained by summing over the other variable:
\[
p_X(x) = \sum_{y} p_{XY}(x,y)
\]
\[
p_Y(y) = \sum_{x} p_{XY}(x,y)
\]

\textbf{Key Idea:} The marginal PMF of one variable is obtained by summing the joint PMF over all possible values of the other variable.

\subsection{Marginal PMF from a Joint PMF Table}

From a joint PMF table:
\begin{itemize}
\item Row sums give $p_X(x)$
\item Column sums give $p_Y(y)$
\end{itemize}

\[
p_X(x_i) = \sum_j p_{ij}, \quad p_Y(y_j) = \sum_i p_{ij}
\]

Also,
\[
\sum_i p_X(x_i) = 1, \quad \sum_j p_Y(y_j) = 1
\]

\subsection{Solved Example}

Given joint PMF:

\[
\begin{array}{c|cc}
 & y=0 & y=1 \\
\hline
x=0 & \frac{1}{8} & \frac{1}{4} \\
x=1 & \frac{3}{8} & \frac{1}{4}
\end{array}
\]

\textbf{Step 1: Compute $p_X(x)$}
\[
p_X(0) = \frac{1}{8} + \frac{1}{4} = \frac{3}{8}
\]
\[
p_X(1) = \frac{3}{8} + \frac{1}{4} = \frac{5}{8}
\]

\textbf{Step 2: Compute $p_Y(y)$}
\[
p_Y(0) = \frac{1}{8} + \frac{3}{8} = \frac{1}{2}
\]
\[
p_Y(1) = \frac{1}{4} + \frac{1}{4} = \frac{1}{2}
\]

\textbf{Final Result:}
\[
p_X(x) =
\begin{cases}
\frac{3}{8}, & x=0 \\
\frac{5}{8}, & x=1
\end{cases}
\quad
p_Y(y) =
\begin{cases}
\frac{1}{2}, & y=0 \\
\frac{1}{2}, & y=1
\end{cases}
\]

\section{Correlation}

\subsection{Definition}

A basic way to measure joint behavior is through the product moment:
\[
R_{XY} = E[XY]
\]

It represents the average value of the product $XY$.

It is a raw correlation measure.

\subsection{Properties and Interpretation}

\[
R_{XY} = E[XY]
\]

\begin{itemize}
\item Depends on the scale of $X$ and $Y$
\item Computed about the origin, not about the means
\item Mixes:
\begin{itemize}
\item Mean values
\item Spread
\item Dependence
\end{itemize}
\end{itemize}

If:
\[
R_{XY} = 0
\]
$\Rightarrow$ Variables are orthogonal about the origin.

\textbf{Transition:} To remove the effect of means, variables are centered $\rightarrow$ leads to covariance.

\section{Covariance}

\subsection{Definition}

\[
\text{Cov}(X,Y) = E[(X - \mu_X)(Y - \mu_Y)]
\]

where:
\[
\mu_X = E[X], \quad \mu_Y = E[Y]
\]

\subsection{Interpretation}

\begin{itemize}
\item Measures how centered variables vary together
\item Indicates direction of joint variation:
\begin{itemize}
\item Positive covariance $\rightarrow$ same direction
\item Negative covariance $\rightarrow$ opposite direction
\item Zero covariance $\rightarrow$ no linear co-variation
\end{itemize}
\item Important: Magnitude depends on units $\rightarrow$ cannot measure strength directly
\end{itemize}

\subsection{Algebraic Form}

Starting with:
\[
\text{Cov}(X,Y) = E[(X - \mu_X)(Y - \mu_Y)]
\]

\[
= E[XY - X\mu_Y - \mu_X Y + \mu_X \mu_Y]
\]

\[
= E[XY] - \mu_Y E[X] - \mu_X E[Y] + \mu_X \mu_Y
\]

Since:
\[
\mu_X = E[X], \quad \mu_Y = E[Y]
\]

\[
\text{Cov}(X,Y) = E[XY] - E[X]E[Y]
\]

\subsection{Properties}

\begin{itemize}
\item Symmetry:
\[
\text{Cov}(X,Y) = \text{Cov}(Y,X)
\]

\item Variance as Special Case:
\[
\text{Cov}(X,X) = \text{Var}(X)
\]

\item Linearity:
\[
\text{Cov}(aX + b, cY + d) = ac\,\text{Cov}(X,Y)
\]

\item Independence:
\[
X \perp Y \Rightarrow \text{Cov}(X,Y) = 0
\]
\end{itemize}

\textbf{Note:} Zero covariance $\Rightarrow$ uncorrelated, but does not imply independence.

\subsection{Example 1}

Given:
\[
Y = -2X + 3
\]

Find $\text{Cov}(X,Y)$.

\textbf{Step 1:}
\[
\text{Cov}(X,Y) = E[XY] - E[X]E[Y]
\]

\textbf{Step 2:}
\[
XY = X(-2X + 3) = -2X^2 + 3X
\]
\[
E[Y] = E[-2X + 3] = -2E[X] + 3
\]

\textbf{Step 3:}
\[
\text{Cov}(X,Y) = E[-2X^2 + 3X] - E[X](-2E[X] + 3)
\]
\[
= -2E[X^2] + 3E[X] + 2(E[X])^2 - 3E[X]
\]
\[
= -2E[X^2] + 2(E[X])^2
\]

\textbf{Step 4:}

Using:
\[
\text{Var}(X) = E[X^2] - (E[X])^2
\]

\[
\text{Cov}(X,Y) = -2\,\text{Var}(X)
\]

\subsection{Pearson’s Correlation Coefficient}

\[
\rho_{XY} = \frac{\text{Cov}(X,Y)}{\sqrt{\text{Var}(X)\text{Var}(Y)}} = \frac{\text{Cov}(X,Y)}{\sigma_X \sigma_Y}
\]

\begin{itemize}
\item Measures strength and direction of linear relationship
\item Normalized version of covariance
\item Unit-free
\end{itemize}

\[
-1 \le \rho_{XY} \le 1
\]

\[
\rho_{XY} > 0: \text{ positive linear trend}
\]
\[
\rho_{XY} < 0: \text{ negative linear trend}
\]
\[
\rho_{XY} = 0: \text{ no linear correlation}
\]

\subsection{Example 2}

Given:
\[
f_{XY}(x,y) =
\begin{cases}
\frac{xy}{96}, & 0 < x < 4, \; 1 < y < 5 \\
0, & \text{otherwise}
\end{cases}
\]

Results (as derived in slides):

\[
E[X] = \frac{8}{3}
\]
\[
E[Y] = \frac{31}{9}
\]

Since $X$ and $Y$ are independent:
\[
E[XY] = E[X]E[Y] = \frac{248}{27}
\]

\[
E[2X + 3Y] = \frac{47}{3}
\]

\[
\text{Var}(X) = \frac{8}{9}
\]
\[
\text{Var}(Y) = \frac{92}{81}
\]

\[
\text{Cov}(X,Y) = 0
\]
\[
\rho_{XY} = 0
\]

\section{Jointly Gaussian Random Variables}

\subsection{Definition}

A pair $(X,Y)$ is jointly Gaussian if:
\begin{itemize}
\item The joint density has the bivariate Gaussian form
\item The marginals are Gaussian
\end{itemize}

\subsection{Applications}

\begin{itemize}
\item Signal processing
\item Communication systems
\item Modeling continuous noisy measurements
\end{itemize}

\subsection{Joint Gaussian PDF}

\[
f_{XY}(x,y) =
\frac{1}{2\pi \sigma_X \sigma_Y \sqrt{1-\rho^2}}
\exp \left(
-\frac{1}{2(1-\rho^2)}
\left[
\left(\frac{x-\mu_X}{\sigma_X}\right)^2
+
\left(\frac{y-\mu_Y}{\sigma_Y}\right)^2
-
2\rho
\left(\frac{x-\mu_X}{\sigma_X}\right)
\left(\frac{y-\mu_Y}{\sigma_Y}\right)
\right]
\right)
\]

Where:
\begin{itemize}
\item $\mu_X, \mu_Y$: means
\item $\sigma_X, \sigma_Y$: standard deviations
\item $\rho$: Pearson correlation coefficient
\end{itemize}

\subsection{Interpretation}

Example (from slide):

\begin{itemize}
\item $X$: CO$_2$ concentration
\item $Y$: temperature anomaly
\end{itemize}

If:
\[
\rho > 0
\]
$\Rightarrow$ Higher CO$_2$ levels tend to occur with higher temperatures.

\end{document} 